% --- IMPORTS ---
\documentclass[leqno]{article}
\usepackage{graphicx}
\usepackage{dsfont}
\usepackage{mathtools}
\usepackage{amssymb}
\usepackage{amsmath}
\usepackage{amsthm}
\usepackage{float}
\usepackage{ragged2e}
\usepackage{tensor}
\usepackage{fancyhdr}
\usepackage{empheq}
\usepackage[most]{tcolorbox}
\usepackage{enumitem}
\usepackage{dirtytalk}
\usepackage{marginnote}
\usepackage{microtype}
\usepackage[
  a4paper,
  left=0.8in,
  right=1in,
  top=1in,
  bottom=0.5in,
  headheight=14pt,
  headsep=0.4in
]{geometry}
\setlength{\parindent}{0cm}
% --- TikZ Packages ---
\usepackage{pgfplots}
\pgfplotsset{compat=1.18}
\usepgfplotslibrary{fillbetween, polar}
\usetikzlibrary{calc, positioning}
\usetikzlibrary{angles, quotes}
\usetikzlibrary{arrows.meta}
\usetikzlibrary{decorations.pathreplacing, decorations.markings}
\usetikzlibrary{patterns}
\usetikzlibrary{intersections}
% --- URL Packages ---
\usepackage[colorlinks=true,linkcolor=red,citecolor=magenta,urlcolor=black]{hyperref}%
\urlstyle{same}
% --- END OF IMPORTS ---

% --- CUSTOM COMMANDS ---
\RenewDocumentCommand{\marks}{o m}{%
  \IfValueTF{#1}%
    {\marginnote{\raggedright\textbf{[#2]}}[#1]}%
    {\marginnote{\raggedright\textbf{[#2]}}}%
}
\newcommand{\vspacerule}[2]{%
  \par\vspace*{#1}%
  \noindent\hrulefill\par
  \vspace*{#2}%
}
\definecolor{scarlet}{RGB}{205, 40, 75}
\newcommand{\questionitem}{%
  \item\phantomsection
  \addcontentsline{toc}{section}{Question \number\value{enumi}}%
}
\renewcommand{\contentsname}{Questions}
% --- END OF CUSTOM COMMANDS ---

% --- HEADER CONFIG ---
\pagestyle{fancy}
\fancyhead{}
\fancyfoot{}
\renewcommand{\headrulewidth}{0.9pt}
\lhead{\textit{Solutions} - Vectors - Reflections in Planes}
\chead{}
\rhead{\href{https://danieljsmith.org}{danieljsmith.org}}
% --- END OF HEADER CONFIG ---

\begin{document}
\vspace*{20pt}
\tableofcontents
\newpage
\begin{enumerate}[
  label=\textbf{\arabic*.},
  leftmargin=*,
  topsep=0pt,
  partopsep=0pt,
  parsep=0pt
]

% ---- Question 1 ----
\questionitem

The line $\ell_1$ is the intersection of the planes
\[
x+y+z=1
\] \\
and
\[
2x-y+z=4
\] \\
The plane $\Pi$ has equation
\[
x+y+2z=3
\] \\
The line $\ell_2$ is the reflection of the line $\ell_1$ in the plane $\Pi$. \\

Find a vector equation of the line $\ell_2$. \marks{7}

\vspace*{30pt}

\begin{tcolorbox}[colback=white, colframe=scarlet, title={\textbf{Solution}}, breakable, grow to left by=6mm, grow to right by=10mm]
Write the reflecting plane in the form
\[
\Pi:\ \mathbf r\cdot \mathbf n=d
\]
so here
\[
\mathbf n=\begin{pmatrix}1\\1\\2\end{pmatrix},
\qquad d=3
\]

To find the reflection of the line \(\ell_1\), we use the reflection formulae for a point and a direction vector in a plane.

First find a direction vector for \(\ell_1\). Since \(\ell_1\) is the intersection of the planes
\[
x+y+z=1
\qquad\text{and}\qquad
2x-y+z=4
\]
their normal vectors are
\[
\mathbf n_1=\begin{pmatrix}1\\1\\1\end{pmatrix},
\qquad
\mathbf n_2=\begin{pmatrix}2\\-1\\1\end{pmatrix}
\]

So a direction vector of \(\ell_1\) is
\begin{align*}
\mathbf d_1&=\mathbf n_1\times \mathbf n_2\\
&=
\begin{vmatrix}
\mathbf i&\mathbf j&\mathbf k\\
1&1&1\\
2&-1&1
\end{vmatrix}\\
&=2\mathbf i+\mathbf j-3\mathbf k\\
&=\begin{pmatrix}2\\1\\-3\end{pmatrix}
\end{align*}

Now find a convenient point on \(\ell_1\). The best choice is the point where \(\ell_1\) meets \(\Pi\), since a point on the mirror plane is unchanged by reflection.

So solve
\begin{align*}
x+y+z&=1\\
2x-y+z&=4\\
x+y+2z&=3
\end{align*}

Subtracting the first equation from the third gives
\begin{align*}
z&=2
\end{align*}

Substituting \(z=2\) into the first two equations:
\begin{align*}
x+y+2&=1\\
2x-y+2&=4
\end{align*}
so
\begin{align*}
x+y&=-1\\
2x-y&=2
\end{align*}

Adding these equations:
\begin{align*}
3x&=1\\
x&=\frac13
\end{align*}

Then
\begin{align*}
y&=-1-\frac13=-\frac43
\end{align*}

Hence the point is
\[
\mathbf q=\begin{pmatrix}\frac13\\-\frac43\\2\end{pmatrix}
\]

Using the point reflection formula
\[
\mathbf q'
=
\mathbf q
-2\,\frac{\mathbf q\cdot\mathbf n-d}{\mathbf n\cdot\mathbf n}\,\mathbf n
\]
we note that \(\mathbf q\) lies on \(\Pi\), so \(\mathbf q\cdot \mathbf n=d\). Therefore
\[
\mathbf q'=\mathbf q
\]

So the reflected line \(\ell_2\) passes through the same point
\[
\begin{pmatrix}\frac13\\-\frac43\\2\end{pmatrix}
\]

Now reflect the direction vector using
\[
\mathbf d_2
=
\mathbf d_1
-2\,\frac{\mathbf d_1\cdot\mathbf n}{\mathbf n\cdot\mathbf n}\,\mathbf n
\]

Compute the dot products:
\begin{align*}
\mathbf d_1\cdot \mathbf n&=
\begin{pmatrix}2\\1\\-3\end{pmatrix}\cdot
\begin{pmatrix}1\\1\\2\end{pmatrix}
=2+1-6=-3\\
\mathbf n\cdot \mathbf n&=
\begin{pmatrix}1\\1\\2\end{pmatrix}\cdot
\begin{pmatrix}1\\1\\2\end{pmatrix}
=1+1+4=6
\end{align*}

Hence
\begin{align*}
\mathbf d_2
&=
\begin{pmatrix}2\\1\\-3\end{pmatrix}
-2\left(\frac{-3}{6}\right)
\begin{pmatrix}1\\1\\2\end{pmatrix}\\
&=
\begin{pmatrix}2\\1\\-3\end{pmatrix}
+\begin{pmatrix}1\\1\\2\end{pmatrix}\\
&=
\begin{pmatrix}3\\2\\-1\end{pmatrix}
\end{align*}

Therefore a vector equation of the reflected line \(\ell_2\) is
\[
\mathbf r=
\begin{pmatrix}\frac13\\-\frac43\\2\end{pmatrix}
+t\begin{pmatrix}3\\2\\-1\end{pmatrix},
\qquad t\in\mathbb{R}
\]

This is a valid vector equation of the line \(\ell_2\).
\end{tcolorbox}


\newpage

% ---- Question 2 ----
\questionitem

The plane $\Pi$ has vector equation
\[
\mathbf{r}=\begin{pmatrix}1\\2\\-1\end{pmatrix}+s\begin{pmatrix}2\\-1\\1\end{pmatrix}+t\begin{pmatrix}1\\1\\0\end{pmatrix}
\] \\
The line $\ell_1$ passes through the point $P(4,1,2)$ and is perpendicular to $\Pi$. \\
The line $\ell_1$ meets $\Pi$ at the point $Q$. \\
\begin{enumerate}[label=\textbf{(\alph*)}]
  \item Find the coordinates of $Q$. \marks{4} \\
\end{enumerate}
Given that the point $R(3,1,0)$ lies on $\Pi$, \\
\begin{enumerate}[label=\textbf{(\alph*)}, resume]
  \item find a Cartesian equation of the plane containing $P$, $Q$ and $R$. \marks{4} \\
\end{enumerate}
The line $\ell_2$ passes through $P$ and $R$. \\
The line $\ell_3$ is the reflection of $\ell_2$ in $\Pi$. \\
\begin{enumerate}[label=\textbf{(\alph*)}, resume]
  \item Find a vector equation of $\ell_3$. \marks{4} \\
\end{enumerate}

\vspace*{30pt}

\begin{tcolorbox}[colback=white, colframe=scarlet, title={\textbf{Solution}}, breakable, grow to left by=6mm, grow to right by=10mm]
\begin{enumerate}[label=\textbf{(\alph*)}]
  \item The plane \(\Pi\) has direction vectors
\[
\begin{pmatrix}2\\-1\\1\end{pmatrix}
\quad \text{and} \quad
\begin{pmatrix}1\\1\\0\end{pmatrix}
\]

So a normal vector to \(\Pi\) is
\begin{align*}
\mathbf n
&=
\begin{pmatrix}2\\-1\\1\end{pmatrix}
\times
\begin{pmatrix}1\\1\\0\end{pmatrix}\\
&=
\begin{pmatrix}
(-1)(0)-1(1)\\
1(1)-2(0)\\
2(1)-(-1)(1)
\end{pmatrix}\\
&=
\begin{pmatrix}-1\\1\\3\end{pmatrix}
\end{align*}

Using the point \(\begin{pmatrix}1\\2\\-1\end{pmatrix}\) on \(\Pi\),
\begin{align*}
d
&=
\begin{pmatrix}1\\2\\-1\end{pmatrix}\cdot
\begin{pmatrix}-1\\1\\3\end{pmatrix}\\
&=-1+2-3\\
&=-2
\end{align*}

Hence \(\Pi\) is
\[
\mathbf r\cdot \begin{pmatrix}-1\\1\\3\end{pmatrix}=-2
\]
so in Cartesian form,
\[
-x+y+3z+2=0
\]

Since \(\ell_1\) is perpendicular to \(\Pi\), it has direction vector \(\mathbf n\). Therefore
\[
\ell_1:\ \mathbf r=
\begin{pmatrix}4\\1\\2\end{pmatrix}
+\lambda
\begin{pmatrix}-1\\1\\3\end{pmatrix}
\]

Using the line--plane formula
\[
\lambda_*=\frac{d-\mathbf a\cdot\mathbf n}{\mathbf b\cdot\mathbf n}
\]
with \(\mathbf a=\begin{pmatrix}4\\1\\2\end{pmatrix}\), \(\mathbf b=\begin{pmatrix}-1\\1\\3\end{pmatrix}\), \(d=-2\),
\begin{align*}
\lambda
&=
\frac{-2-\begin{pmatrix}4\\1\\2\end{pmatrix}\cdot\begin{pmatrix}-1\\1\\3\end{pmatrix}}
{\begin{pmatrix}-1\\1\\3\end{pmatrix}\cdot\begin{pmatrix}-1\\1\\3\end{pmatrix}}\\
&=
\frac{-2-(-4+1+6)}{1+1+9}\\
&=
\frac{-2-3}{11}\\
&=
-\frac{5}{11}
\end{align*}

So
\begin{align*}
Q
&=
\begin{pmatrix}4\\1\\2\end{pmatrix}
-\frac{5}{11}
\begin{pmatrix}-1\\1\\3\end{pmatrix}\\
&=
\begin{pmatrix}
4+\frac{5}{11}\\[2pt]
1-\frac{5}{11}\\[2pt]
2-\frac{15}{11}
\end{pmatrix}\\
&=
\begin{pmatrix}
\frac{49}{11}\\[2pt]
\frac{6}{11}\\[2pt]
\frac{7}{11}
\end{pmatrix}
\end{align*}

Hence
\[
Q\left(\frac{49}{11},\frac{6}{11},\frac{7}{11}\right)
\]

  \item A plane containing \(P,Q,R\) contains the vectors \(\overrightarrow{PQ}\) and \(\overrightarrow{QR}\).

First,
\begin{align*}
\overrightarrow{PQ}
&=
\begin{pmatrix}\frac{49}{11}\\[2pt]\frac{6}{11}\\[2pt]\frac{7}{11}\end{pmatrix}
-
\begin{pmatrix}4\\1\\2\end{pmatrix}\\
&=
\begin{pmatrix}
\frac{5}{11}\\[2pt]
-\frac{5}{11}\\[2pt]
-\frac{15}{11}
\end{pmatrix}
=
\frac{5}{11}
\begin{pmatrix}
1\\-1\\-3
\end{pmatrix}
\end{align*}

Also,
\begin{align*}
\overrightarrow{QR}
&=
\begin{pmatrix}3\\1\\0\end{pmatrix}
-
\begin{pmatrix}\frac{49}{11}\\[2pt]\frac{6}{11}\\[2pt]\frac{7}{11}\end{pmatrix}\\
&=
\begin{pmatrix}
-\frac{16}{11}\\[2pt]
\frac{5}{11}\\[2pt]
-\frac{7}{11}
\end{pmatrix}
=
\frac{1}{11}
\begin{pmatrix}
-16\\5\\-7
\end{pmatrix}
\end{align*}

So a normal vector to the plane \(PQR\) is
\begin{align*}
\mathbf n_1
&=
\begin{pmatrix}1\\-1\\-3\end{pmatrix}
\times
\begin{pmatrix}-16\\5\\-7\end{pmatrix}\\
&=
\begin{pmatrix}
(-1)(-7)-(-3)(5)\\
(-3)(-16)-1(-7)\\
1(5)-(-1)(-16)
\end{pmatrix}\\
&=
\begin{pmatrix}
22\\55\\-11
\end{pmatrix}\\
&=
11\begin{pmatrix}2\\5\\-1\end{pmatrix}
\end{align*}

Hence we may take the normal vector as
\[
\begin{pmatrix}2\\5\\-1\end{pmatrix}
\]

Using the point \(P(4,1,2)\), the plane equation is
\[
\begin{pmatrix}2\\5\\-1\end{pmatrix}\cdot
\left(
\begin{pmatrix}x\\y\\z\end{pmatrix}
-
\begin{pmatrix}4\\1\\2\end{pmatrix}
\right)=0
\]

So
\begin{align*}
2(x-4)+5(y-1)-(z-2)&=0\\
2x-8+5y-5-z+2&=0\\
2x+5y-z-11&=0
\end{align*}

Therefore a Cartesian equation of the plane containing \(P,Q,R\) is
\[
2x+5y-z-11=0
\]

  \item The line \(\ell_2\) passes through \(R(3,1,0)\) and \(P(4,1,2)\), so
\[
\ell_2:\ \mathbf r=
\begin{pmatrix}3\\1\\0\end{pmatrix}
+\mu
\begin{pmatrix}1\\0\\2\end{pmatrix}
\]

Write \(\Pi\) in the form \(\mathbf r\cdot\mathbf n=d\), where
\[
\mathbf n=\begin{pmatrix}-1\\1\\3\end{pmatrix},
\qquad
d=-2
\]

Since \(R\) lies on \(\Pi\), it is unchanged by reflection. Indeed,
\[
\begin{pmatrix}3\\1\\0\end{pmatrix}\cdot
\begin{pmatrix}-1\\1\\3\end{pmatrix}
=-3+1+0=-2=d
\]

So to find \(\ell_3\), we reflect the direction vector
\[
\mathbf b=\begin{pmatrix}1\\0\\2\end{pmatrix}
\]

Using the reflection formula for a direction vector,
\[
\mathbf b'
=\mathbf b-2\,\frac{\mathbf b\cdot\mathbf n}{\mathbf n\cdot\mathbf n}\,\mathbf n
\]

Now
\begin{align*}
\mathbf b\cdot\mathbf n
&=
\begin{pmatrix}1\\0\\2\end{pmatrix}\cdot
\begin{pmatrix}-1\\1\\3\end{pmatrix}\\
&=-1+0+6\\
&=5
\end{align*}

and
\begin{align*}
\mathbf n\cdot\mathbf n
&=
\begin{pmatrix}-1\\1\\3\end{pmatrix}\cdot
\begin{pmatrix}-1\\1\\3\end{pmatrix}\\
&=1+1+9\\
&=11
\end{align*}

Therefore
\begin{align*}
\mathbf b'
&=
\begin{pmatrix}1\\0\\2\end{pmatrix}
-2\cdot \frac{5}{11}
\begin{pmatrix}-1\\1\\3\end{pmatrix}\\
&=
\begin{pmatrix}1\\0\\2\end{pmatrix}
-\begin{pmatrix}-\frac{10}{11}\\[2pt]\frac{10}{11}\\[2pt]\frac{30}{11}\end{pmatrix}\\
&=
\begin{pmatrix}
\frac{21}{11}\\[2pt]
-\frac{10}{11}\\[2pt]
-\frac{8}{11}
\end{pmatrix}
\end{align*}

This is parallel to \(\begin{pmatrix}21\\-10\\-8\end{pmatrix}\). Hence
\[
\ell_3:\ \mathbf r=
\begin{pmatrix}3\\1\\0\end{pmatrix}
+\mu
\begin{pmatrix}21\\-10\\-8\end{pmatrix}
\]

So a vector equation of \(\ell_3\) is
\[
\mathbf r=
\begin{pmatrix}3\\1\\0\end{pmatrix}
+\mu
\begin{pmatrix}21\\-10\\-8\end{pmatrix}
\]
\end{enumerate}
\end{tcolorbox}


\newpage

% ---- Question 3 ----
\questionitem

A plane passes through the points $C(1,-1,2)$, $D(3,0,1)$ and $E(2,2,2)$. \\

The point $A$ has coordinates $(4,-2,7)$ and the point $B$ is the reflection of $A$ in the plane. \\

Find the coordinates of the point $B$. \marks{7}

\vspace*{30pt}

\begin{tcolorbox}[colback=white, colframe=scarlet, title={\textbf{Solution}}, breakable, grow to left by=6mm, grow to right by=10mm]
Let the plane be \(\Pi\).

Two direction vectors in the plane are
\[
\overrightarrow{CD}=(3-1,\;0-(-1),\;1-2)=(2,1,-1)
\]
and
\[
\overrightarrow{CE}=(2-1,\;2-(-1),\;2-2)=(1,3,0)
\]

So a normal vector to the plane is
\begin{align*}
\mathbf n
&=\overrightarrow{CD}\times\overrightarrow{CE}\\
&=
\begin{vmatrix}
\mathbf i & \mathbf j & \mathbf k\\
2 & 1 & -1\\
1 & 3 & 0
\end{vmatrix}\\
&=\mathbf i(1\cdot 0-(-1)\cdot 3)-\mathbf j(2\cdot 0-(-1)\cdot 1)+\mathbf k(2\cdot 3-1\cdot 1)\\
&=3\mathbf i-\mathbf j+5\mathbf k\\
&=(3,-1,5)
\end{align*}

Using the form \(\mathbf r\cdot\mathbf n=d\), we find \(d\) from point \(C(1,-1,2)\):
\begin{align*}
d&=(1,-1,2)\cdot(3,-1,5)\\
&=1\cdot 3+(-1)\cdot(-1)+2\cdot 5\\
&=3+1+10\\
&=14
\end{align*}

Hence the plane is
\[
\mathbf r\cdot(3,-1,5)=14
\]
so in Cartesian form,
\[
3x-y+5z-14=0
\]

Now use the reflection formula for a point in a plane:
\[
\mathbf p'
=\mathbf p
-2\,\frac{\mathbf p\cdot\mathbf n-d}{\mathbf n\cdot\mathbf n}\,\mathbf n
\]

Here
\[
\mathbf p=(4,-2,7),\qquad \mathbf n=(3,-1,5),\qquad d=14
\]

First calculate \(\mathbf p\cdot\mathbf n-d\):
\begin{align*}
\mathbf p\cdot\mathbf n-d
&=(4,-2,7)\cdot(3,-1,5)-14\\
&=4\cdot 3+(-2)\cdot(-1)+7\cdot 5-14\\
&=12+2+35-14\\
&=35
\end{align*}

Also,
\[
\mathbf n\cdot\mathbf n=3^2+(-1)^2+5^2=9+1+25=35
\]

Substituting into the reflection formula:
\begin{align*}
\mathbf p'
&=(4,-2,7)-2\left(\frac{35}{35}\right)(3,-1,5)\\
&=(4,-2,7)-2(3,-1,5)\\
&=(4,-2,7)-(6,-2,10)\\
&=(-2,0,-3)
\end{align*}

Therefore the coordinates of \(B\) are
\[
B=(-2,0,-3)
\]

As a quick check, the closest point on the plane to \(A\) is
\[
\mathbf p-\frac{\mathbf p\cdot\mathbf n-d}{\mathbf n\cdot\mathbf n}\mathbf n
=(4,-2,7)-\frac{35}{35}(3,-1,5)
=(1,-1,2)=C
\]
so \(C\) is the midpoint of \(AB\), which confirms the result.

Final answer: \(B=(-2,0,-3)\).
\end{tcolorbox}


\newpage

% ---- Question 4 ----
\questionitem

A line $\ell$ has equation
\[
\mathbf{r}=\begin{pmatrix}4\\2\\-3\end{pmatrix}+\lambda\begin{pmatrix}2\\-1\\1\end{pmatrix}
\] \\
Another line $m$ has equation
\[
\mathbf{r}=\begin{pmatrix}1\\1\\2\end{pmatrix}+\mu\begin{pmatrix}2\\1\\0\end{pmatrix}
\] \\
A plane $\Pi$ contains $m$ and is parallel to $\ell$. \\
\begin{enumerate}[label=\textbf{(\roman*)}]
  \item Find an equation of $\Pi$, giving your answer in the form $\mathbf{r}\cdot\mathbf{n}=p$. \marks{4} \\
  \item Find the distance between $\ell$ and $\Pi$. \marks{4} \\
  \item Find an equation of the line which is the reflection of $\ell$ in $\Pi$, giving your answer in the form $\mathbf{r}=\mathbf{a}+t\mathbf{b}$. \marks{4} \\
\end{enumerate}

\vspace*{30pt}

\begin{tcolorbox}[colback=white, colframe=scarlet, title={\textbf{Solution}}, breakable, grow to left by=6mm, grow to right by=10mm]
\begin{enumerate}[label=\textbf{(\roman*)}]
  \item Let the direction vector of line \(\ell\) be
\[
\mathbf b=\begin{pmatrix}2\\-1\\1\end{pmatrix}
\]
and the direction vector of line \(m\) be
\[
\mathbf d=\begin{pmatrix}2\\1\\0\end{pmatrix}
\]

Since \(\Pi\) contains \(m\) and is parallel to \(\ell\), both \(\mathbf d\) and \(\mathbf b\) lie in the plane, so a normal vector is
\begin{align*}
\mathbf n&=\mathbf d\times \mathbf b \\
&=\begin{pmatrix}2\\1\\0\end{pmatrix}\times \begin{pmatrix}2\\-1\\1\end{pmatrix} \\
&=\begin{pmatrix}
1\cdot 1-0\cdot (-1)\\
-\left(2\cdot 1-0\cdot 2\right)\\
2\cdot (-1)-1\cdot 2
\end{pmatrix} \\
&=\begin{pmatrix}1\\-2\\-4\end{pmatrix}
\end{align*}

A point on \(m\) is \(\begin{pmatrix}1\\1\\2\end{pmatrix}\), so for the plane equation \(\mathbf r\cdot \mathbf n=p\),
\begin{align*}
p&=\begin{pmatrix}1\\-2\\-4\end{pmatrix}\cdot \begin{pmatrix}1\\1\\2\end{pmatrix} \\
&=1-2-8 \\
&=-9
\end{align*}

Hence the plane is
\[
\mathbf r\cdot \begin{pmatrix}1\\-2\\-4\end{pmatrix}=-9
\]

So an equation of \(\Pi\) is
\[
\mathbf r\cdot \begin{pmatrix}1\\-2\\-4\end{pmatrix}=-9
\]

  \item From the line--plane formula, if \(\mathbf b\cdot \mathbf n=0\), then the line is parallel to the plane.

Here
\begin{align*}
\mathbf b\cdot \mathbf n&=\begin{pmatrix}2\\-1\\1\end{pmatrix}\cdot \begin{pmatrix}1\\-2\\-4\end{pmatrix} \\
&=2+2-4 \\
&=0
\end{align*}

So \(\ell\parallel \Pi\).

Take the point \(P=(4,2,-3)\) on \(\ell\). Using the point-to-plane distance formula
\[
\mathrm{dist}=\frac{|\mathbf p\cdot \mathbf n-d|}{|\mathbf n|}
\]
with \(\mathbf p=\begin{pmatrix}4\\2\\-3\end{pmatrix}\), \(\mathbf n=\begin{pmatrix}1\\-2\\-4\end{pmatrix}\) and \(d=-9\),

\begin{align*}
\mathbf p\cdot \mathbf n&=\begin{pmatrix}4\\2\\-3\end{pmatrix}\cdot \begin{pmatrix}1\\-2\\-4\end{pmatrix} \\
&=4-4+12 \\
&=12
\end{align*}

Also,
\begin{align*}
|\mathbf n|&=\sqrt{1^2+(-2)^2+(-4)^2} \\
&=\sqrt{1+4+16} \\
&=\sqrt{21}
\end{align*}

Therefore
\begin{align*}
\mathrm{dist}&=\frac{|12-(-9)|}{\sqrt{21}} \\
&=\frac{21}{\sqrt{21}} \\
&=\sqrt{21}
\end{align*}

So the distance between \(\ell\) and \(\Pi\) is
\[
\sqrt{21}
\]

  \item Use the reflection formulae for a plane \(\Pi:\mathbf r\cdot \mathbf n=d\).

For a point,
\[
\mathbf p'
=\mathbf p
-2\,\frac{\mathbf p\cdot\mathbf n-d}{\mathbf n\cdot\mathbf n}\,\mathbf n
\]

For a direction vector,
\[
\mathbf b'
=\mathbf b
-2\,\frac{\mathbf b\cdot\mathbf n}{\mathbf n\cdot\mathbf n}\,\mathbf n
\]

Here
\[
\mathbf n=\begin{pmatrix}1\\-2\\-4\end{pmatrix},\qquad d=-9
\]

First reflect the direction vector of \(\ell\):
\begin{align*}
\mathbf b\cdot \mathbf n&=0
\end{align*}
so
\begin{align*}
\mathbf b'
&=\mathbf b-2\,\frac{0}{\mathbf n\cdot\mathbf n}\,\mathbf n \\
&=\mathbf b \\
&=\begin{pmatrix}2\\-1\\1\end{pmatrix}
\end{align*}

So the direction is unchanged.

Now reflect the point \(P=(4,2,-3)\) on \(\ell\):
\begin{align*}
\mathbf n\cdot \mathbf n&=1^2+(-2)^2+(-4)^2=21 \\
\mathbf p\cdot \mathbf n-d&=12-(-9)=21
\end{align*}

Hence
\begin{align*}
\mathbf p'
&=\begin{pmatrix}4\\2\\-3\end{pmatrix}
-2\left(\frac{21}{21}\right)\begin{pmatrix}1\\-2\\-4\end{pmatrix} \\
&=\begin{pmatrix}4\\2\\-3\end{pmatrix}
-2\begin{pmatrix}1\\-2\\-4\end{pmatrix} \\
&=\begin{pmatrix}4\\2\\-3\end{pmatrix}
-\begin{pmatrix}2\\-4\\-8\end{pmatrix} \\
&=\begin{pmatrix}2\\6\\5\end{pmatrix}
\end{align*}

Therefore the reflected line is
\[
\mathbf r=\begin{pmatrix}2\\6\\5\end{pmatrix}+t\begin{pmatrix}2\\-1\\1\end{pmatrix}
\]

So an equation of the reflected line is
\[
\mathbf r=\begin{pmatrix}2\\6\\5\end{pmatrix}+t\begin{pmatrix}2\\-1\\1\end{pmatrix}
\]
\end{enumerate}
\end{tcolorbox}


\newpage

% ---- Question 5 ----
\questionitem

The line $\ell_1$ has vector equation
\[
\mathbf{r}=\begin{pmatrix}1\\4\\0\end{pmatrix}+\lambda\begin{pmatrix}2\\1\\1\end{pmatrix}
\] \\
The plane $\Pi_1$ passes through the point $(2,0,0)$ and has normal vector
\[
\begin{pmatrix}3\\-1\\2\end{pmatrix}
\] \\
\begin{enumerate}[label=\textbf{(\alph*)}]
  \item Find the point of intersection of $\ell_1$ and $\Pi_1$. \marks{2} \\
\end{enumerate}
The line $\ell_2$ is the reflection of the line $\ell_1$ in the plane $\Pi_1$. \\
\begin{enumerate}[label=\textbf{(\alph*)}, resume]
  \item Show that a vector equation for the line $\ell_2$ is
  \[
  \mathbf{r}=\begin{pmatrix}4\\3\\2\end{pmatrix}+\mu\begin{pmatrix}-1\\2\\-1\end{pmatrix}
  \]
  where $\mu$ is a scalar parameter. \marks{5}\\
\end{enumerate}
The plane $\Pi_2$ contains the line $\ell_1$ and the line $\ell_2$. \\
\begin{enumerate}[label=\textbf{(\alph*)}, resume]
  \item Determine a vector equation for the line of intersection of $\Pi_1$ and $\Pi_2$. \marks{2} \\
\end{enumerate}

\vspace*{30pt}

\begin{tcolorbox}[colback=white, colframe=scarlet, title={\textbf{Solution}}, breakable, grow to left by=6mm, grow to right by=10mm]
\begin{enumerate}[label=\textbf{(\alph*)}]
  \item The plane $\Pi_1$ passes through $(2,0,0)$ and has normal vector $\begin{pmatrix}3\\-1\\2\end{pmatrix}$, so its equation is
\[
\begin{pmatrix}3\\-1\\2\end{pmatrix}\cdot
\left(
\begin{pmatrix}x\\y\\z\end{pmatrix}
-
\begin{pmatrix}2\\0\\0\end{pmatrix}
\right)=0
\]
Hence
\[
3(x-2)-y+2z=0
\]
so
\[
3x-y+2z=6
\]

A general point on $\ell_1$ is
\[
(x,y,z)=(1+2\lambda,\ 4+\lambda,\ \lambda)
\]

Substituting into the plane equation:
\begin{align*}
3(1+2\lambda)-(4+\lambda)+2\lambda&=6\\
3+6\lambda-4-\lambda+2\lambda&=6\\
7\lambda-1&=6\\
7\lambda&=7\\
\lambda&=1
\end{align*}

So the point of intersection is
\[
(1+2,\ 4+1,\ 1)=(3,5,1)
\]

Therefore $\ell_1$ meets $\Pi_1$ at $(3,5,1)$.

  \item Let
\[
\mathbf a=\begin{pmatrix}1\\4\\0\end{pmatrix},
\qquad
\mathbf u=\begin{pmatrix}2\\1\\1\end{pmatrix},
\qquad
\mathbf n=\begin{pmatrix}3\\-1\\2\end{pmatrix}
\]
so that
\[
\ell_1:\ \mathbf r=\mathbf a+\lambda \mathbf u
\]

Since $\Pi_1$ passes through $(2,0,0)$,
\[
\Pi_1:\ \mathbf r\cdot \mathbf n=6
\]
because
\[
\begin{pmatrix}2\\0\\0\end{pmatrix}\cdot \begin{pmatrix}3\\-1\\2\end{pmatrix}=6
\]

Using the reflection formulae for a plane $\mathbf r\cdot \mathbf n=6$,

\[
\mathbf p'
=
\mathbf p
-2\,\frac{\mathbf p\cdot\mathbf n-6}{\mathbf n\cdot\mathbf n}\,\mathbf n
\]
for a point, and
\[
\mathbf u'
=
\mathbf u
-2\,\frac{\mathbf u\cdot\mathbf n}{\mathbf n\cdot\mathbf n}\,\mathbf n
\]
for a direction vector.

First reflect the point $\mathbf a$.

\begin{align*}
\mathbf a\cdot\mathbf n
&=
\begin{pmatrix}1\\4\\0\end{pmatrix}\cdot
\begin{pmatrix}3\\-1\\2\end{pmatrix}\\
&=3-4+0\\
&=-1
\end{align*}

Also,
\[
\mathbf n\cdot\mathbf n=3^2+(-1)^2+2^2=14
\]

So
\begin{align*}
\mathbf a'
&=
\mathbf a
-2\,\frac{\mathbf a\cdot\mathbf n-6}{\mathbf n\cdot\mathbf n}\,\mathbf n\\
&=
\mathbf a
-2\,\frac{-1-6}{14}\,\mathbf n\\
&=
\mathbf a
-2\left(\frac{-7}{14}\right)\mathbf n\\
&=
\mathbf a+\mathbf n\\
&=
\begin{pmatrix}1\\4\\0\end{pmatrix}
+
\begin{pmatrix}3\\-1\\2\end{pmatrix}\\
&=
\begin{pmatrix}4\\3\\2\end{pmatrix}
\end{align*}

Now reflect the direction vector $\mathbf u$.

\begin{align*}
\mathbf u\cdot\mathbf n
&=
\begin{pmatrix}2\\1\\1\end{pmatrix}\cdot
\begin{pmatrix}3\\-1\\2\end{pmatrix}\\
&=6-1+2\\
&=7
\end{align*}

Hence
\begin{align*}
\mathbf u'
&=
\mathbf u
-2\,\frac{\mathbf u\cdot\mathbf n}{\mathbf n\cdot\mathbf n}\,\mathbf n\\
&=
\mathbf u
-2\,\frac{7}{14}\,\mathbf n\\
&=
\mathbf u-\mathbf n\\
&=
\begin{pmatrix}2\\1\\1\end{pmatrix}
-
\begin{pmatrix}3\\-1\\2\end{pmatrix}\\
&=
\begin{pmatrix}-1\\2\\-1\end{pmatrix}
\end{align*}

Therefore the reflected line $\ell_2$ has vector equation
\[
\mathbf r=
\begin{pmatrix}4\\3\\2\end{pmatrix}
+
\mu
\begin{pmatrix}-1\\2\\-1\end{pmatrix}
\]

which is the required result.

  \item From part (a), the point
\[
P=(3,5,1)
\]
lies on $\Pi_1$ and on $\ell_1$.

Since $\ell_1\subset \Pi_2$, the point $P$ also lies on $\Pi_2$. So the line of intersection of $\Pi_1$ and $\Pi_2$ passes through $P$.

A direction vector of $\ell_1$ is
\[
\mathbf u_1=\begin{pmatrix}2\\1\\1\end{pmatrix}
\]
and a direction vector of $\ell_2$ is
\[
\mathbf u_2=\begin{pmatrix}-1\\2\\-1\end{pmatrix}
\]

Because $\Pi_2$ contains both lines, it contains both direction vectors, so it also contains
\[
\mathbf u_1+\mathbf u_2
=
\begin{pmatrix}2\\1\\1\end{pmatrix}
+
\begin{pmatrix}-1\\2\\-1\end{pmatrix}
=
\begin{pmatrix}1\\3\\0\end{pmatrix}
\]

To check that this is also parallel to $\Pi_1$, take the dot product with the normal to $\Pi_1$:
\[
\begin{pmatrix}1\\3\\0\end{pmatrix}\cdot
\begin{pmatrix}3\\-1\\2\end{pmatrix}
=3-3+0=0
\]

So $\begin{pmatrix}1\\3\\0\end{pmatrix}$ is a direction vector for the line $\Pi_1\cap\Pi_2$.

Hence a vector equation for the line of intersection is
\[
\mathbf r=
\begin{pmatrix}3\\5\\1\end{pmatrix}
+t\begin{pmatrix}1\\3\\0\end{pmatrix}
\]
where $t$ is a scalar parameter.
\end{enumerate}
\end{tcolorbox}


\newpage

% ---- Question 6 ----
\questionitem

The triangle $ABC$ lies in the plane $\Pi$. The vertices are
\[
A(1,0,0),\ B(2,2,1) \text{ and } C(1,1,c)
\] \\
where $c$ is a constant. \\
\begin{enumerate}[label=\textbf{(\alph*)}]
  \item Find, in terms of $c$, $\overrightarrow{BA}\times\overrightarrow{BC}$. \marks{3} \\
\end{enumerate}
Given that $\overrightarrow{BA}\times\overrightarrow{BC}=d\mathbf{i}+2\mathbf{j}-\mathbf{k}$, where $d$ is a constant, \\
\begin{enumerate}[label=\textbf{(\alph*)}, resume]
  \item find the value of $c$ and show that $d=-3$. \marks{2} \\
  \item find an equation of $\Pi$ in the form $\mathbf{r}\cdot\mathbf{n}=p$, where $p$ is a constant. \marks{3} \\
\end{enumerate}
The point $S$ has position vector $3\mathbf{i}+\mathbf{j}+2\mathbf{k}$. The point $S'$ is the image of $S$ under reflection in $\Pi$. \\
\begin{enumerate}[label=\textbf{(\alph*)}, resume]
  \item Find the position vector of $S'$. \marks{5} \\
\end{enumerate}

\vspace*{30pt}

\begin{tcolorbox}[colback=white, colframe=scarlet, title={\textbf{Solution}}, breakable, grow to left by=6mm, grow to right by=10mm]
\begin{enumerate}[label=\textbf{(\alph*)}]
  \item First find the two side vectors from \(B\):
\[
\overrightarrow{BA}=\overrightarrow{OA}-\overrightarrow{OB}=(1,0,0)-(2,2,1)=(-1,-2,-1)
\]
\[
\overrightarrow{BC}=\overrightarrow{OC}-\overrightarrow{OB}=(1,1,c)-(2,2,1)=(-1,-1,c-1)
\]

Now take the cross product:
\begin{align*}
\overrightarrow{BA}\times \overrightarrow{BC}
&=
\begin{vmatrix}
\mathbf i & \mathbf j & \mathbf k\\
-1 & -2 & -1\\
-1 & -1 & c-1
\end{vmatrix} \\
&=\mathbf i\bigl((-2)(c-1)-(-1)(-1)\bigr)
-\mathbf j\bigl((-1)(c-1)-(-1)(-1)\bigr)
+\mathbf k\bigl((-1)(-1)-(-2)(-1)\bigr) \\
&=\mathbf i\bigl(-2c+2-1\bigr)
-\mathbf j\bigl(-c+1-1\bigr)
+\mathbf k(1-2) \\
&=(1-2c)\mathbf i+c\mathbf j-\mathbf k
\end{align*}

So,
\[
\overrightarrow{BA}\times \overrightarrow{BC}=(1-2c)\mathbf i+c\mathbf j-\mathbf k
\]

  \item We are given that
\[
\overrightarrow{BA}\times \overrightarrow{BC}=d\mathbf i+2\mathbf j-\mathbf k
\]

From part (a),
\[
\overrightarrow{BA}\times \overrightarrow{BC}=(1-2c)\mathbf i+c\mathbf j-\mathbf k
\]

Comparing coefficients of \(\mathbf j\),
\[
c=2
\]

Then comparing coefficients of \(\mathbf i\),
\[
d=1-2c=1-2(2)=1-4=-3
\]

Hence,
\[
c=2 \quad \text{and} \quad d=-3
\]

  \item A normal vector to the plane \(\Pi\) is
\[
\mathbf n=\overrightarrow{BA}\times \overrightarrow{BC}
\]

Using \(c=2\), part (a) gives
\[
\mathbf n=(1-2(2),2,-1)=(-3,2,-1)
\]

So we can take
\[
\mathbf n=-3\mathbf i+2\mathbf j-\mathbf k
\]

The plane has equation \(\mathbf r\cdot \mathbf n=p\). Since point \(A(1,0,0)\) lies on \(\Pi\),
\begin{align*}
p&=\overrightarrow{OA}\cdot \mathbf n \\
&=(1,0,0)\cdot(-3,2,-1) \\
&=-3
\end{align*}

Therefore the equation of \(\Pi\) is
\[
\mathbf r\cdot(-3\mathbf i+2\mathbf j-\mathbf k)=-3
\]

Equivalently, in Cartesian form,
\[
-3x+2y-z=-3
\]

  \item Let the position vector of \(S\) be
\[
\mathbf p=3\mathbf i+\mathbf j+2\mathbf k=(3,1,2)
\]

From part (c), the plane is
\[
\mathbf r\cdot \mathbf n=d
\]
with
\[
\mathbf n=(-3,2,-1), \qquad d=-3
\]

Using the reflection formula for a point in a plane,
\[
\mathbf p'
=\mathbf p
-2\,\frac{\mathbf p\cdot\mathbf n-d}
{\mathbf n\cdot\mathbf n}\,\mathbf n
\]

We now calculate the required quantities:
\begin{align*}
\mathbf p\cdot\mathbf n
&=(3,1,2)\cdot(-3,2,-1) \\
&=3(-3)+1(2)+2(-1) \\
&=-9+2-2 \\
&=-9
\end{align*}

\begin{align*}
\mathbf n\cdot\mathbf n
&=(-3)^2+2^2+(-1)^2 \\
&=9+4+1 \\
&=14
\end{align*}

So
\begin{align*}
\mathbf p'
&=\mathbf p-2\frac{-9-(-3)}{14}\mathbf n \\
&=\mathbf p-2\frac{-6}{14}\mathbf n \\
&=\mathbf p+\frac{12}{14}\mathbf n \\
&=\mathbf p+\frac67\mathbf n
\end{align*}

Substituting \(\mathbf p=(3,1,2)\) and \(\mathbf n=(-3,2,-1)\),
\begin{align*}
\mathbf p'
&=(3,1,2)+\frac67(-3,2,-1) \\
&=\left(3-\frac{18}{7},\,1+\frac{12}{7},\,2-\frac67\right) \\
&=\left(\frac{21}{7}-\frac{18}{7},\,\frac{7}{7}+\frac{12}{7},\,\frac{14}{7}-\frac67\right) \\
&=\left(\frac37,\frac{19}{7},\frac87\right)
\end{align*}

Hence the position vector of \(S'\) is
\[
\frac37\mathbf i+\frac{19}{7}\mathbf j+\frac87\mathbf k
\]
\end{enumerate}
\end{tcolorbox}


\newpage

% ---- Question 7 ----
\questionitem

The line $\ell$ has equation
\[
\mathbf{r}=(\mathbf{i}+\mathbf{k})+\lambda(\mathbf{i}+2\mathbf{j}+3\mathbf{k})
\] \\
where $\lambda$ is a scalar parameter, and the plane $\Pi$ has equation
\[
\mathbf{r}\cdot(2\mathbf{i}-\mathbf{j}+\mathbf{k})=9
\] \\
\begin{enumerate}[label=\textbf{(\alph*)}]
  \item Find the coordinates of the point of intersection of $\ell$ and $\Pi$. \marks{4} \\
\end{enumerate}
The perpendicular to $\Pi$ from the point $A\,(0,4,1)$ meets $\Pi$ at the point $B$. \\
\begin{enumerate}[label=\textbf{(\alph*)}, resume]
  \item Verify that the coordinates of $B$ are $(4,2,3)$. \marks{3} \\
\end{enumerate}
The point $A\,(0,4,1)$ is reflected in the plane $\Pi$ to give the image point $A'$. \\
\begin{enumerate}[label=\textbf{(\alph*)}, resume]
  \item Find the coordinates of the point $A'$. \marks{2} \\
  \item Find an equation for the line obtained by reflecting the line $\ell$ in the plane $\Pi$, giving your answer in the form
  \[
  \mathbf{r}\times\mathbf{a}=\mathbf{b}
  \]
  \marks[-20pt]{4}
\end{enumerate}

\vspace*{30pt}

\begin{tcolorbox}[colback=white, colframe=scarlet, title={\textbf{Solution}}, breakable, grow to left by=6mm, grow to right by=10mm]
\begin{enumerate}[label=\textbf{(\alph*)}]
  \item Write the line \(\ell\) in component form:
  \[
  \mathbf r=(1,0,1)+\lambda(1,2,3)
  \]
  so
  \[
  (x,y,z)=(1+\lambda,\,2\lambda,\,1+3\lambda)
  \]

  The plane \(\Pi\) has equation
  \[
  \mathbf r\cdot(2,-1,1)=9
  \]
  so in Cartesian form,
  \[
  2x-y+z=9
  \]

  Substituting the line coordinates into the plane equation:
  \begin{align*}
  2(1+\lambda)-2\lambda+(1+3\lambda)&=9\\
  2+2\lambda-2\lambda+1+3\lambda&=9\\
  3+3\lambda&=9\\
  3\lambda&=6\\
  \lambda&=2
  \end{align*}

  Hence the point of intersection is
  \[
  (x,y,z)=(1+2,\,2\cdot 2,\,1+3\cdot 2)=(3,4,7)
  \]

  Therefore, the line \(\ell\) meets the plane \(\Pi\) at \((3,4,7)\).

  \item A normal vector to the plane \(\Pi\) is
  \[
  \mathbf n=(2,-1,1)
  \]

  So the perpendicular through \(A(0,4,1)\) has equation
  \[
  \mathbf r=(0,4,1)+t(2,-1,1)
  \]

  Hence points on this line have coordinates
  \[
  (x,y,z)=(2t,\,4-t,\,1+t)
  \]

  Substituting into the plane equation \(2x-y+z=9\):
  \begin{align*}
  2(2t)-(4-t)+(1+t)&=9\\
  4t-4+t+1+t&=9\\
  6t-3&=9\\
  6t&=12\\
  t&=2
  \end{align*}

  So
  \[
  B=(2\cdot 2,\,4-2,\,1+2)=(4,2,3)
  \]

  Hence the coordinates of \(B\) are verified to be \((4,2,3)\).

  \item Using the reflection formula for a point in the plane \(\Pi:\mathbf r\cdot\mathbf n=d\),
  \[
  \mathbf p'
  =\mathbf p-2\,\frac{\mathbf p\cdot\mathbf n-d}{\mathbf n\cdot\mathbf n}\,\mathbf n
  \]

  Here
  \[
  \mathbf p=(0,4,1),\qquad \mathbf n=(2,-1,1),\qquad d=9
  \]

  First calculate:
  \[
  \mathbf p\cdot\mathbf n=0\cdot 2+4(-1)+1\cdot 1=-3,
  \qquad
  \mathbf n\cdot\mathbf n=2^2+(-1)^2+1^2=6
  \]

  Therefore
  \begin{align*}
  \mathbf p'
  &=(0,4,1)-2\,\frac{-3-9}{6}(2,-1,1)\\
  &=(0,4,1)-2\left(\frac{-12}{6}\right)(2,-1,1)\\
  &=(0,4,1)+4(2,-1,1)\\
  &=(8,0,5)
  \end{align*}

  So the image point is \(A'=(8,0,5)\).

  \item Let \(P\) be the point where \(\ell\) meets \(\Pi\). From part (a),
  \[
  P=(3,4,7)
  \]

  Since \(P\) lies in the mirror plane, it stays fixed under reflection.

  The direction vector of \(\ell\) is
  \[
  \mathbf d=(1,2,3)
  \]

  Using the reflection formula for a direction vector,
  \[
  \mathbf d'
  =\mathbf d-2\,\frac{\mathbf d\cdot\mathbf n}{\mathbf n\cdot\mathbf n}\,\mathbf n
  \]
  with \(\mathbf n=(2,-1,1)\).

  First calculate:
  \[
  \mathbf d\cdot\mathbf n=1\cdot 2+2(-1)+3\cdot 1=3,
  \qquad
  \mathbf n\cdot\mathbf n=6
  \]

  So
  \begin{align*}
  \mathbf d'
  &=(1,2,3)-2\left(\frac{3}{6}\right)(2,-1,1)\\
  &=(1,2,3)-(2,-1,1)\\
  &=(-1,3,2)
  \end{align*}

  Hence the reflected line \(\ell'\) is
  \[
  \mathbf r=(3,4,7)+\mu(-1,3,2)
  \]

  To write this in the form \(\mathbf r\times\mathbf a=\mathbf b\), take
  \[
  \mathbf a=(-1,3,2)
  \]

  Then
  \[
  \mathbf r=(3,4,7)+\mu\mathbf a
  \]
  implies
  \[
  \mathbf r\times\mathbf a=((3,4,7)+\mu\mathbf a)\times\mathbf a=(3,4,7)\times\mathbf a
  \]
  since \(\mathbf a\times\mathbf a=\mathbf 0\).

  Now calculate
  \begin{align*}
  (3,4,7)\times(-1,3,2)
  &=\begin{vmatrix}
  \mathbf i & \mathbf j & \mathbf k\\
  3&4&7\\
  -1&3&2
  \end{vmatrix}\\
  &=\mathbf i(4\cdot 2-7\cdot 3)-\mathbf j(3\cdot 2-7(-1))+\mathbf k(3\cdot 3-4(-1))\\
  &=\mathbf i(8-21)-\mathbf j(6+7)+\mathbf k(9+4)\\
  &=-13\mathbf i-13\mathbf j+13\mathbf k
  \end{align*}

  Therefore an equation of the reflected line is
  \[
  \mathbf r\times(-\mathbf i+3\mathbf j+2\mathbf k)=-13\mathbf i-13\mathbf j+13\mathbf k
  \]

  So
  \[
  \mathbf a=(-1,3,2),\qquad \mathbf b=(-13,-13,13)
  \]
\end{enumerate}
\end{tcolorbox}


\end{enumerate}
\end{document}
